\section*{Tools Used}
\addcontentsline{toc}{section}{Tools Used}
\label{sec:tools}

\paragraph{AI assistance.} Claude (Anthropic) was used to help write the project's code and to draft its documents, including the README, the tutorial and this report. The four members reviewed what it produced, ran the tests and the evaluation, and are responsible for the content. The evaluation numbers in this report come from the results files written by the code (Tables~\ref{tab:retrieval}--\ref{tab:reasons} and Figure~\ref{fig:sweep} are generated from them by \code{scripts/make_report_assets.py}), not from the assistant. \TODO{each member: confirm or complete this statement before signing}

\paragraph{Models used by the system.} The language models that answer the questions and judge faithfulness are free models served by OpenRouter (\code{nvidia/nemotron-3-ultra-550b-a55b:free} and \code{google/gemma-4-31b-it:free}). The embedding model is \code{BAAI/bge-m3}; the gate cross-encoder is \path{cross-encoder/ms-marco-MiniLM-L6-v2}; \path{ncbi/MedCPT-Cross-Encoder} and \path{BAAI/bge-reranker-v2-m3} were compared with it.

\paragraph{Libraries.} PyTorch (CPU), sentence-transformers, chromadb, pypdf, cryptography, Gradio and pytest; the report uses \LaTeX{} with PGFPlots and TikZ.

\paragraph{Borrowed code and tutorials.} BM25 follows two tutorials listed by \emph{build-your-own-x} \citep{buildyourownx}: term and document frequencies after \citet{loria_tfidf}, and the inverted index after \citet{boyter_vsm}; the formula is from \citet{robertson2009bm25}, with the idf variant of Lucene \citep{lucenebm25}. The function \code{rrf} has the structure of \code{reciprocal_rank_fusion} in the \emph{rag-from-scratch} tutorial \citep{ragfromscratch}. Both are credited in comments in the code.
